%----------
%	MODELS - LLMs
%----------	

In contrast to the implementations of ML and DL models described in Sections~\ref{machine_learning_models} and ~\ref{implementation_dl}, and as seen in Figure~\ref{fig:llms_models_implementation}, no pre-processing techniques were applied to the text, except for filtering by language as discussed in Section~\ref{languages}. This decision was made because generative models are highly effective at processing natural language without the need for extensive pre-processing. Additionally, since we are employing a zero-shot learning approach, as detailed in Section~\ref{zero_shot}, we did not perform any data splitting or balancing techniques, as the entire dataset was classified.


\begin{figure}[H]
	\ffigbox[\FBwidth]
	{\caption{Diagram of the Implementation of GenAI Models}\label{fig:llms_models_implementation}} 
 % [Aquí ponemos como queremos que aparezca en el índice, sin la referencia]{Aquí como queremos que aparezca en el documento, con la referencia}
	{\includegraphics[scale=0.45]{Plantilla_TFG_ingles_2019/imagenes/TFG Pipeline-LLMs.drawio.pdf}}
\end{figure}


The first step in the implementation of Generative Models was Prompt Engineering, as seen in Figure~\ref{fig:prompt_engineering_implementation}, where we designed prompts that guide the model to generate the desired output. This step included adding context, specific instructions, categories and class definitions, as described in Section~\ref{target_variables_defs}. In particular, we experimented with two types of contexts: "Context 0", which was essentially empty, and "Context 1", which contained the detailed information described in Section~\ref{contextual_information}. Additionally, we developed four variations of prompts. These variations included two suffixes, "0" and "1", to indicate the presence or absence of class descriptions for each task, and two language options, "EN" and "ES", to differentiate whether the prompt information was presented in English or Spanish. For the responses, the prompts instructed the model to make predictions using the following format: '0' and '1' for "Análisis General", '0' through '3' for "Contenido Negativo", and '0' through '2' for "Insultos". A collection of these contexts and prompts can be found in Appendices~\ref{appendix_zeroshot_context}-\ref{appendix_zeroshot_prompts_insultos}.


\begin{figure}[H]
	\ffigbox[\FBwidth]
	{\caption{Diagram of the Implementation of Prompt Engineering for GenAI Models}\label{fig:prompt_engineering_implementation}} 
 % [Aquí ponemos como queremos que aparezca en el índice, sin la referencia]{Aquí como queremos que aparezca en el documento, con la referencia}
	{\includegraphics[scale=0.33]{Plantilla_TFG_ingles_2019/imagenes/TFG Pipeline-GenAI.drawio.pdf}}
\end{figure}


After setting up the prompts, we moved on to implementing the models using LM Studio\footnote{From \url{https://lmstudio.ai/}, last visited 16/07/24}, a tool that facilitates the deployment of large language models locally. Specifically, we utilized the models detailed in Section~\ref{zero_shot}, and tuned the "Temperature" hyperparameter (0.1 and 0.5), totalling 64 experiments per task. To effectively use LM Studio, we started a local server that acted as an endpoint for making HTTP requests. We leveraged the "Multi-Model Session" feature, which provides access to "JSON Mode", allowing us to constrain each model's output format to our needs. In this case, we decided to format the response as \textit{\{"prediction": ""\}}. This setup enabled us to send tweets along with the prompts via the LM Studio API to the server and receive predictions from the selected model in a structured manner. The processing of the responses was straightforward as they only needed to be parsed and cast to numbers, in order to be able to evaluate the predictions in different metrics. Finally, a summary of the experiments and the different metrics were logged into an Excel file called \texttt{"Training\_GenAI.xlsx"} (Appendix~\ref{appendix_code}).

As explained before, the number of experiments was equal across all classes, so we would expect the time spent on each task to be proportional to the number of instances per task, Consequently, more time was spent on "Análisis General" than on "Contenido Negativo" and "Insultos". Therefore, Figure~\ref{fig:genai_models_total_performance} details the distribution of time allocated for each classification task, with execution time sorted according to the task with the highest number of instances.


\begin{figure}[H]
	\ffigbox[\FBwidth]
	{\caption{Total Performance Time (in hours) per Classification Task for GenAI Experiments}\label{fig:genai_models_total_performance}} 
 % [Aquí ponemos como queremos que aparezca en el índice, sin la referencia]{Aquí como queremos que aparezca en el documento, con la referencia}
	{\includegraphics[scale=0.23]{Plantilla_TFG_ingles_2019/imagenes/GenAI_total_performance_time.pdf}}
\end{figure}
